\section{Geometric Interfaces Used by the Proof}
\label{sec:reader-framework}

% This section records the exact outputs of the structural and local calculations.
% Their derivations are postponed.  The distinction is deliberate: the body uses
% only the statements below, while the appendices verify them on explicit
% real-variable domains.

\subsection{Structural reduction}

\begin{proposition}[Reader structural reduction]
\label{prop:body-structural-reduction}
Every hypothetical cover has the seven distinct roles described in
Section~\ref{sec:introduction}.  The center has exactly one of the types
$\CEzero,\CEone,\CEtwo$, every vertex role has exactly one of the types
$\Vdzero,\Vdone,\Vdtwo,\Tlike$, and the actual reaches determine exactly one
entry of Table~\ref{tab:routing}.  In the CE1/CE2 case the center contains exactly
one radial midpoint and
\[
 q=N_++m.
\]
If the vertex roles alone cover the boundary, strict handoffs can be selected
while preserving whether the actual supercritical pattern has none, exactly
one, or at least two indices.
\end{proposition}

\begin{proof}
Distinctness is Lemma~\ref{lem:distinct-roles}; the center and vertex
classifications are Propositions~\ref{prop:ce-classification} and
\ref{prop:vertex-classification}; the unique midpoint is
Proposition~\ref{prop:unique-center-midpoint}; the short-role identity is
Lemma~\ref{lem:short-role-count}; strict selection is
Proposition~\ref{prop:strict-handoffs}; and the disjoint routing assertion is
Proposition~\ref{prop:exhaustive-structural-reduction}.  Their geometric proofs
are collected in Appendix~\ref{sec:structural-reductions}.
\end{proof}

\subsection{The local admissible set}

\begin{definition}[Actual and selected skeleton-V-triangle reach coordinates]
\label{def:body-vtriangle-coordinates}
All indices lie in $\mathbb Z/6\mathbb Z$. For a tagged open or closed
$V_i$-triangle, the distinguished vertex $V_i$ is part of the role data and
the reach functions are given by the following three suprema. At the closed
vertex role $T_i$, define
\[
 \begin{aligned}
 A(T_i)&=
 \sup_{\substack{t\in[0,1]\\
 V_i+t(V_{i-1}-V_i)\in T_i}}t,\\
 B(T_i)&=
 \sup_{\substack{t\in[0,1]\\
 V_i+t(V_{i+1}-V_i)\in T_i}}t,\\
 C(T_i)&=
 \sup_{\substack{t\in[0,1]\\
 V_i+t(O-V_i)\in T_i}}t.
 \end{aligned}
\]
The functions give the incoming reach toward $V_{i-1}$, the outgoing reach
toward $V_{i+1}$, and the own-radial reach toward $O$, respectively. Since
$T_i=\overline{U_i}$, they satisfy
\[
 A(U_i)=A(T_i),\qquad B(U_i)=B(T_i),\qquad C(U_i)=C(T_i).
\]
For compact indexed formulas, put
\[
 (A_i,B_i,C_i):=(A(T_i),B(T_i),C(T_i)).
\]
The triple $(A(T_i),B(T_i),C(T_i))$ is the \emph{actual V-triangle reach data}. A
\emph{selected demand triple} is any $(a_i,b_i,c_i)\in[0,1]^3$ satisfying
\[
 (a_i,b_i,c_i)\le(A(T_i),B(T_i),C(T_i))
\]
coordinatewise.  Define
\[
 x_i=1-a_i,\qquad y_i=b_i.
\]
The actual V triangle is nonsupercritical, critical, or strictly supercritical
according as
\[
 A(T_i)+B(T_i)<1,\qquad
 A(T_i)+B(T_i)=1,\qquad
 A(T_i)+B(T_i)>1,
\]
respectively.  In particular,
\[
 N_+=\left\lvert
 \left\lbrace i\in\mathbb Z/6\mathbb Z:A(T_i)+B(T_i)>1\right\rbrace
 \right\rvert
\]
is defined from actual data, never from a selected triple.
\end{definition}

\begin{definition}[Local demand hull and admissible set]
\label{def:body-local-admissible-set}
Let
\[
 u=\left(\frac12,\frac{\sqrt3}{2}\right),\qquad
 v=\left(\frac12,-\frac{\sqrt3}{2}\right).
\]
For $(a,b,c)\in[0,1]^3$, define
\[
 K(a,b,c)=\conv\{0,au,bv,c(u+v)\}
\]
\[
 q_\theta=(\cos\theta,\sin\theta),\qquad
 r_\theta=
 \left(\cos\left(\theta+\frac{\pi}{3}\right),
       \sin\left(\theta+\frac{\pi}{3}\right)\right),
\]
\[
 \mathfrak E_1=
 \left\{
 \conv\{z,z+q_\theta,z+r_\theta\}:
 z\in\mathbb R^2,\ 0\le\theta<2\pi
 \right\}.
\]
Define
\[
 \mathcal A=
 \left\{(a,b,c)\in[0,1]^3:
 \exists\,\Delta\in\mathfrak E_1,\quad K(a,b,c)\subseteq\Delta
 \right\}.
\]
Thus a closed vertex role with actual V-triangle reach data $(A,B,C)$ realizes every
$(a,b,c)\in\mathcal A$ satisfying
\[
 0\le a\le A,\qquad 0\le b\le B,\qquad 0\le c\le C.
\]
\end{definition}

\begin{figure}[H]
\centering
\resizebox{.98\linewidth}{!}{%
\begin{tikzpicture}[x=1cm,y=1cm,>=Latex]
  \node[panellabel,anchor=west] at (.05,2.62)
    {Local demand data inside the regular-hexagon skeleton};

  \begin{scope}[shift={(2.30,.25)},scale=2.10]
    \HexCoords
    \DrawHexagon
    \DrawRays

    \coordinate (ain) at ($(V0)!.36!(V5)$);
    \coordinate (Ain) at ($(V0)!.56!(V5)$);
    \coordinate (bout) at ($(V0)!.31!(V1)$);
    \coordinate (Bout) at ($(V0)!.52!(V1)$);
    \coordinate (crad) at ($(V0)!.43!(O)$);
    \coordinate (Crad) at ($(V0)!.65!(O)$);

    \draw[actualreach] (V0)--(Ain);
    \draw[demand] (V0)--(ain);
    \draw[actualreach] (V0)--(Bout);
    \draw[demand] (V0)--(bout);
    \draw[actualreach] (V0)--(Crad);
    \draw[radialdemand] (V0)--(crad);

    \fill[black] (V0) circle (.8pt);
    \node[figlabel,right] at (V0) {$V_i$};
    \node[figlabel,above] at ($(V0)!.68!(V1)$)
      {$e_{i,i+1}$};
    \node[figlabel,below] at ($(V0)!.70!(V5)$)
      {$e_{i-1,i}$};
    \node[figlabel,above right,text=DemandRed] at ($(V0)!.17!(V1)$)
      {$b_i$};
    \node[figlabel,above left,text=VertexOrange] at ($(V0)!.46!(V1)$)
      {$B_i$};
    \node[figlabel,below right,text=DemandRed] at ($(V0)!.17!(V5)$)
      {$a_i$};
    \node[figlabel,below left,text=VertexOrange] at ($(V0)!.48!(V5)$)
      {$A_i$};
    \node[figlabel,above,text=RadialTeal] at ($(V0)!.30!(O)$)
      {$c_i$};
    \node[figlabel,below,text=VertexOrange] at ($(V0)!.57!(O)$)
      {$C_i$};
    \node[figlabel,atlasbox,align=left,anchor=west] at (-1.05,-1.48)
      {$0\le a_i\le A_i,\quad0\le b_i\le B_i,\quad
        0\le c_i\le C_i$\\
       $x_i=1-a_i,\qquad y_i=b_i$};
  \end{scope}

  \begin{scope}[shift={(8.15,.20)},scale=2.30]
    \coordinate (Q0) at (0,0);
    \coordinate (Qa) at (.34,.589);
    \coordinate (Qb) at (.48,-.831);
    \coordinate (Qc) at (.67,0);
    \draw[->,black!45] (-.05,0)--(.95,0)
      node[right,figlabel] {$u+v$};
    \draw[->,black!45] (Q0)--(.54,.935)
      node[above,figlabel] {$u$};
    \draw[->,black!45] (Q0)--(.54,-.935)
      node[below,figlabel] {$v$};
    \draw[thick,fill=DemandRed!8]
      (Q0)--(Qa)--(Qc)--(Qb)--cycle;
    \fill (Q0) circle (.7pt) node[left,figlabel] {$0$};
    \fill (Qa) circle (.7pt) node[above left,figlabel] {$au$};
    \fill (Qb) circle (.7pt) node[below left,figlabel] {$bv$};
    \fill (Qc) circle (.7pt) node[above right,figlabel] {$c(u+v)$};
    \node[figlabel,atlasbox,align=center] at (.43,-1.25)
      {$K(a,b,c)=\conv
       \{0,au,bv,c(u+v)\}$\\
       $(a,b,c)\in\mathcal A\iff\Lambda(K(a,b,c))\le1$};
  \end{scope}

  \draw[flow] (4.80,.25)--node[above,figlabel]
    {translate \(V_i\) to \(0\)} (6.55,.25);
\end{tikzpicture}
}
\caption{Actual and selected V-triangle reach coordinates on the full hexagon skeleton,
together with the rigid local demand hull.  The least equilateral enclosure
side of that hull is the scalar controlling both Strategy~2 transfers and
Strategy~4 witnesses.}
\label{fig:local-demand-hull}
\end{figure}

\begin{figure}[H]
\centering
\resizebox{.98\linewidth}{!}{%
\begin{tikzpicture}[x=1cm,y=1cm,>=Latex]
  \node[panellabel,anchor=west] at (.05,3.72)
    {The oriented unit-triangle family and admissible local hulls};
  \node[figlabel,anchor=east,text=black!60] at (13.45,3.42)
    {a representative member is placed on the full hexagon skeleton};

  \begin{scope}[shift={(2.30,1.30)},scale=1.82]
    \HexCoords
    \DrawHexagon
    \DrawRays

    \coordinate (Z) at (-.42,-.28);
    \coordinate (Zq) at (.46,-.28);
    \coordinate (Zr) at (.02,.49);
    \draw[fill=DemandRed!9,draw=DemandRed,line width=.8pt]
      (Z)--(Zq)--(Zr)--cycle;
    \draw[-{Latex[length=1.2mm]},GapPurple,line width=.7pt]
      (Z)--(Zq);
    \draw[-{Latex[length=1.2mm]},RadialTeal,line width=.7pt]
      (Z)--(Zr);
    \fill[black] (Z) circle (.75pt);
    \node[figlabel,below left] at (Z) {$z$};
    \node[figlabel,below,text=GapPurple] at ($(Z)!.55!(Zq)$)
      {$q_\theta$};
    \node[figlabel,above left,text=RadialTeal] at ($(Z)!.55!(Zr)$)
      {$r_\theta$};
  \end{scope}

  \node[figlabel,atlasbox,align=left,anchor=west] at (4.70,2.62)
    {$\displaystyle
      q_\theta:=(\cos\theta,\sin\theta),\qquad
      r_\theta:=
      \left(\cos\!\left(\theta+\frac\pi3\right),
      \sin\!\left(\theta+\frac\pi3\right)\right)$};

  \node[figlabel,atlasbox,align=left,anchor=west] at (4.70,1.22)
    {$\displaystyle
      \mathfrak E_1:=
      \left\{
      \conv\{z,z+q_\theta,z+r_\theta\}:
      z\in\mathbb R^2,\ 0\le\theta<2\pi
      \right\}$};

  \node[figlabel,atlasbox,align=left,anchor=west] at (4.70,-.18)
    {$\displaystyle
      K(a,b,c):=\conv\{0,au,bv,c(u+v)\},$\\[-1pt]
     $\displaystyle
      \mathcal A:=
      \left\{(a,b,c)\in[0,1]^3:
      \exists\,\Delta\in\mathfrak E_1,
      K(a,b,c)\subseteq\Delta
      \right\}$};
\end{tikzpicture}
}
\caption{The oriented unit-triangle family $\mathfrak E_1$ and the formal
containment definition of the local admissible set $\mathcal A$.}
\label{fig:strategy2-unit-triangle-family-skeleton}
\end{figure}

\begin{definition}[Equilateral enclosure gauge]
\label{def:body-enclosure-gauge}
Define
\[
 \mathsf R:\mathbb R^2\to\mathbb R^2,\qquad
 \mathsf R=
 \begin{pmatrix}
 -1/2&-\sqrt3/2\\
 \sqrt3/2&-1/2
 \end{pmatrix}.
\]
For every nonempty compact $K\subset\mathbb R^2$, define
\[
 h_K:\mathbb S^1\to\mathbb R,\qquad
 h_K(n):=\max_{z\in K}\langle z,n\rangle,
\]
\[
 \Lambda(K):=\frac2{\sqrt3}
 \min_{\lVert n\rVert=1}\sum_{j=0}^2h_K(\mathsf R^jn).
\]
\end{definition}
Proposition~\ref{prop:universal-enclosure-gauge} says that $\Lambda(K)$ is the
least side length of an enclosing closed equilateral triangle.

\begin{proposition}[Geometric structure of the admissible set]
\label{prop:body-admissible-description}
The local admissible set $\mathcal A$ has the following properties.
\begin{enumerate}
\item It is compact, coordinatewise downward closed, and invariant under
interchanging the two boundary coordinates $a$ and $b$.
\item
\[
 (a,b,c)\in\mathcal A
 \quad\Longleftrightarrow\quad
 \Lambda\bigl(K(a,b,c)\bigr)\le1.
\]
\item For $(b,c)$, $(a,c)$, and $(a,b)$ in the corresponding two-coordinate
projections of $\mathcal A$, respectively, define the one-coordinate fibers
\[
 \begin{aligned}
 \mathcal A_{b,c}^{(1)}
 &=\left\{a\ge0:(a,b,c)\in\mathcal A\right\},\\
 \mathcal A_{a,c}^{(2)}
 &=\left\{b\ge0:(a,b,c)\in\mathcal A\right\},\\
 \mathcal A_{a,b}^{(3)}
 &=\left\{c\ge0:(a,b,c)\in\mathcal A\right\}.
 \end{aligned}
\]
Each of these fibers is a nonempty compact initial interval.  In particular,
each maximum used for the Strategy~2 transfers in
Definition~\ref{def:body-transfer-alphabet} is attained.
\end{enumerate}
The exact support expressions, their cell conditions, and the resulting
piecewise equations for $\mathcal A$ are stated and proved in
Appendix~\ref{sec:admissible-set-derivation}.
\end{proposition}

\begin{proof}
Compactness and downward closure follow directly from the containment
definition and compactness of the oriented unit-triangle family modulo
translation.  Reflection in the radial axis exchanges $a$ and $b$.  The gauge
equivalence is Proposition~\ref{prop:universal-enclosure-gauge}.  For each
fixed pair in the corresponding projection, compactness and coordinatewise
downward closure give a nonempty compact initial-interval fiber.  Moreover,
$K(a,0,c)\subseteq\conv\{0,u,u+v\}$ for $a,c\in[0,1]$, so the fibers used for
the later Strategy~2 transfers are nonempty.  The exact cell formulas used in
the plan above are proved in
Appendix~\ref{sec:admissible-set-derivation}.
\end{proof}


Our plan is to split $(a,b,c)$-space into finitely many regions, use one
explicit formula on each region, compare the finitely many candidate enclosing
triangles, and declare $(a,b,c)$ admissible exactly when the smallest candidate
has side at most $1$.


\begin{figure}[H]
\centering
\resizebox{.98\linewidth}{!}{%
\begin{tikzpicture}[x=1cm,y=1cm,>=Latex]
  \node[panellabel,anchor=west] at (.05,3.30)
    {Equilateral enclosure gauge on the hexagon skeleton};
  \node[figlabel,anchor=east,text=black!60] at (13.35,3.02)
    {schematic support directions; formulas and signatures are exact};

  \begin{scope}[shift={(2.35,1.18)},scale=1.82]
    \HexCoords
    \DrawHexagon
    \DrawRays

    \coordinate (K0) at (-.42,-.20);
    \coordinate (K1) at (-.15,.42);
    \coordinate (K2) at (.48,.25);
    \coordinate (K3) at (.37,-.38);
    \draw[fill=DemandRed!10,draw=DemandRed,line width=.75pt]
      (K0)--(K1)--(K2)--(K3)--cycle;
    \node[figlabel,text=DemandRed] at (.02,.02) {$K$};

    \draw[flow] (O)--(.05,.78)
      node[above,figlabel] {$n$};
    \draw[flow] (O)--(-.70,-.38)
      node[below left,figlabel] {$\mathsf R n$};
    \draw[flow] (O)--(.66,-.43)
      node[below right,figlabel] {$\mathsf R^2n$};
    \draw[inactive] (-.58,.48)--(.55,.48);
    \node[figlabel,above] at (0,.48) {$h_K(n)$};
    \fill[CenterBlue] (O) circle (.85pt);
  \end{scope}

  \node[figlabel,atlasbox,align=left,anchor=west] at (4.75,2.42)
    {$\displaystyle
      \mathsf R:\mathbb R^2\to\mathbb R^2,\qquad
      \mathsf R=
      \begin{pmatrix}
       -1/2&-\sqrt3/2\\
       \sqrt3/2&-1/2
      \end{pmatrix}$};

  \node[figlabel,atlasbox,align=left,anchor=west] at (4.75,1.38)
    {For nonempty compact \(K\subset\mathbb R^2\),\\[-2pt]
     \(\displaystyle
       h_K:\mathbb S^1\to\mathbb R,\qquad
       h_K(n):=\max_{z\in K}\langle z,n\rangle.\)};

  \node[figlabel,atlasbox,align=left,anchor=west] at (4.75,.10)
    {$\displaystyle
      \Lambda(K):=\frac2{\sqrt3}
      \min_{\lVert n\rVert=1}
      \sum_{j=0}^{2}h_K(\mathsf R^j n)$\\
     \(\Lambda(K)=\) least side length of a closed equilateral
     triangle containing \(K\).};

  \node[figlabel,atlasbox,align=left,anchor=west] at (4.75,-1.18)
    {$\displaystyle
      (a,b,c)\in\mathcal A
      \iff\Lambda(K(a,b,c))\le1.$};
\end{tikzpicture}
}
\caption{The support-function definition of the equilateral enclosure gauge,
placed in the full skeleton containing the local demand hull.  This figure is
a geometric sketch; the exact caliper equations are in Appendix~\ref{sec:admissible-set-derivation}.}
\label{fig:strategy2-enclosure-gauge-skeleton}
\end{figure}

\subsection{One signed CE1/CE2 interface}

\begin{definition}[Positive part and center-boundary contribution]
\label{def:body-positive-part}
For $t\in\mathbb R$, define
\[
 [t]_+=\max\{t,0\}.
\]
For a closed center role $T_C$, define
\[
 L_{\partial H}(T_C)=\Haus(T_C\cap\partial H).
\]
\end{definition}

\begin{definition}[Geometric signed CE1/CE2 center data]
\label{def:body-signed-center-data}
Normalize one positive-length center trace to $e_{0,1}$.  The signed affine
normal form of Appendix~\ref{sec:signed-ce12-reductions} assigns to the
normalized center triangle a parameter triple $(R,\alpha,\delta)$ in its exact
signed domain $\mathcal D_{\rm sgn}$; write the resulting closed center
triangle as $T_C(R,\alpha,\delta)$.

Let $\mathcal I_{\mathrm R}=T_C\cap e_{0,1}$ be the normalized active trace,
and let $\mathcal I_{\mathrm L}=T_C\cap e_{5,0}$ be the possible companion
trace.  Denote their lengths by $\ell_R$ and $\ell_L$.  The signed companion
coordinate $\Delta_L$ is the affine quantity defined in Appendix
\ref{sec:signed-ce12-reductions}; its sign records whether the companion trace
has positive length.

For each $i\in\mathbb Z/6\mathbb Z$, let $d_i^C$ denote the farthest parameter
reached by $T_C$ on the radial arm $r_i$, measured from $O$ toward $V_i$, and
put $c_i^C=1-d_i^C$.  Thus $d_i^C$ is the center-side radial exit and $c_i^C$
is the complementary radial demand imposed on the V triangle $T_i$.
\end{definition}

\begin{proposition}[Signed center interface]
\label{prop:body-signed-center-interface}
For the normalization of Definition~\ref{def:body-signed-center-data}, the
active trace $\mathcal I_{\mathrm R}$ has positive length.  The center type is
read from the companion trace:
\[
 T_C\text{ is CE1}
 \quad\Longleftrightarrow\quad
 \ell_L=0,
 \qquad
 T_C\text{ is CE2}
 \quad\Longleftrightarrow\quad
 \ell_L>0.
\]
Equivalently, CE2 is characterized by $\Delta_L>0$, while CE1 has
$\Delta_L\le0$; the equality case is a single companion contact point.
Moreover,
\[
 L_{\partial H}(T_C)=\ell_R+\ell_L,
\]
and the six exits $d_i^C$ and complementary demands $c_i^C$ are continuous
piecewise-algebraic functions of the signed parameters.  The bounds needed in
the reader-facing proof are
\[
 L_{\partial H}(T_C)\le \frac{\sqrt3}{2}-\frac34
 \quad\text{in CE1},
 \qquad
 L_{\partial H}(T_C)<\frac12
 \quad\text{in CE2}.
\]
All exact affine side equations, trace endpoints, exit formulas, and sign-cell
conditions are collected in Appendix~\ref{sec:signed-ce12-reductions}.
\end{proposition}

\begin{proof}
This is Proposition~\ref{prop:signed-center-normal-form} and
Corollary~\ref{cor:signed-center-boundary}.  Their exact calculation is kept in
Appendix~\ref{sec:signed-ce12-reductions}; only the geometric classification
and the two boundary bounds are used in the body.
\end{proof}
